% Options for packages loaded elsewhere
\PassOptionsToPackage{unicode}{hyperref}
\PassOptionsToPackage{hyphens}{url}
\PassOptionsToPackage{dvipsnames,svgnames,x11names}{xcolor}
\documentclass[
  10pt,
]{ctexart}
\usepackage{xcolor}
\usepackage[margin=2cm]{geometry}
\usepackage{amsmath,amssymb}
\setcounter{secnumdepth}{-\maxdimen} % remove section numbering
\usepackage{iftex}
\ifPDFTeX
  \usepackage[T1]{fontenc}
  \usepackage[utf8]{inputenc}
  \usepackage{textcomp} % provide euro and other symbols
\else % if luatex or xetex
  \usepackage{unicode-math} % this also loads fontspec
  \defaultfontfeatures{Scale=MatchLowercase}
  \defaultfontfeatures[\rmfamily]{Ligatures=TeX,Scale=1}
\fi
\usepackage{lmodern}
\ifPDFTeX\else
  % xetex/luatex font selection
  \setmainfont[]{DejaVu Sans}
\fi
% Use upquote if available, for straight quotes in verbatim environments
\IfFileExists{upquote.sty}{\usepackage{upquote}}{}
\IfFileExists{microtype.sty}{% use microtype if available
  \usepackage[]{microtype}
  \UseMicrotypeSet[protrusion]{basicmath} % disable protrusion for tt fonts
}{}
\makeatletter
\@ifundefined{KOMAClassName}{% if non-KOMA class
  \IfFileExists{parskip.sty}{%
    \usepackage{parskip}
  }{% else
    \setlength{\parindent}{0pt}
    \setlength{\parskip}{6pt plus 2pt minus 1pt}}
}{% if KOMA class
  \KOMAoptions{parskip=half}}
\makeatother
\usepackage{color}
\usepackage{fancyvrb}
\newcommand{\VerbBar}{|}
\newcommand{\VERB}{\Verb[commandchars=\\\{\}]}
\DefineVerbatimEnvironment{Highlighting}{Verbatim}{commandchars=\\\{\},breaklines=true,breakanywhere=true}
% Add ',fontsize=\small' for more characters per line
\newenvironment{Shaded}{}{}
\newcommand{\AlertTok}[1]{\textcolor[rgb]{1.00,0.00,0.00}{\textbf{#1}}}
\newcommand{\AnnotationTok}[1]{\textcolor[rgb]{0.38,0.63,0.69}{\textbf{\textit{#1}}}}
\newcommand{\AttributeTok}[1]{\textcolor[rgb]{0.49,0.56,0.16}{#1}}
\newcommand{\BaseNTok}[1]{\textcolor[rgb]{0.25,0.63,0.44}{#1}}
\newcommand{\BuiltInTok}[1]{\textcolor[rgb]{0.00,0.50,0.00}{#1}}
\newcommand{\CharTok}[1]{\textcolor[rgb]{0.25,0.44,0.63}{#1}}
\newcommand{\CommentTok}[1]{\textcolor[rgb]{0.38,0.63,0.69}{\textit{#1}}}
\newcommand{\CommentVarTok}[1]{\textcolor[rgb]{0.38,0.63,0.69}{\textbf{\textit{#1}}}}
\newcommand{\ConstantTok}[1]{\textcolor[rgb]{0.53,0.00,0.00}{#1}}
\newcommand{\ControlFlowTok}[1]{\textcolor[rgb]{0.00,0.44,0.13}{\textbf{#1}}}
\newcommand{\DataTypeTok}[1]{\textcolor[rgb]{0.56,0.13,0.00}{#1}}
\newcommand{\DecValTok}[1]{\textcolor[rgb]{0.25,0.63,0.44}{#1}}
\newcommand{\DocumentationTok}[1]{\textcolor[rgb]{0.73,0.13,0.13}{\textit{#1}}}
\newcommand{\ErrorTok}[1]{\textcolor[rgb]{1.00,0.00,0.00}{\textbf{#1}}}
\newcommand{\ExtensionTok}[1]{#1}
\newcommand{\FloatTok}[1]{\textcolor[rgb]{0.25,0.63,0.44}{#1}}
\newcommand{\FunctionTok}[1]{\textcolor[rgb]{0.02,0.16,0.49}{#1}}
\newcommand{\ImportTok}[1]{\textcolor[rgb]{0.00,0.50,0.00}{\textbf{#1}}}
\newcommand{\InformationTok}[1]{\textcolor[rgb]{0.38,0.63,0.69}{\textbf{\textit{#1}}}}
\newcommand{\KeywordTok}[1]{\textcolor[rgb]{0.00,0.44,0.13}{\textbf{#1}}}
\newcommand{\NormalTok}[1]{#1}
\newcommand{\OperatorTok}[1]{\textcolor[rgb]{0.40,0.40,0.40}{#1}}
\newcommand{\OtherTok}[1]{\textcolor[rgb]{0.00,0.44,0.13}{#1}}
\newcommand{\PreprocessorTok}[1]{\textcolor[rgb]{0.74,0.48,0.00}{#1}}
\newcommand{\RegionMarkerTok}[1]{#1}
\newcommand{\SpecialCharTok}[1]{\textcolor[rgb]{0.25,0.44,0.63}{#1}}
\newcommand{\SpecialStringTok}[1]{\textcolor[rgb]{0.73,0.40,0.53}{#1}}
\newcommand{\StringTok}[1]{\textcolor[rgb]{0.25,0.44,0.63}{#1}}
\newcommand{\VariableTok}[1]{\textcolor[rgb]{0.10,0.09,0.49}{#1}}
\newcommand{\VerbatimStringTok}[1]{\textcolor[rgb]{0.25,0.44,0.63}{#1}}
\newcommand{\WarningTok}[1]{\textcolor[rgb]{0.38,0.63,0.69}{\textbf{\textit{#1}}}}
\usepackage{longtable,booktabs,array}
\usepackage{caption}
\captionsetup[table]{skip=6pt}
\newcounter{none} % for unnumbered tables
\usepackage{calc} % for calculating minipage widths
% Correct order of tables after \paragraph or \subparagraph
\usepackage{etoolbox}
\makeatletter
\patchcmd\longtable{\par}{\if@noskipsec\mbox{}\fi\par}{}{}
\makeatother
% Allow footnotes in longtable head/foot
\IfFileExists{footnotehyper.sty}{\usepackage{footnotehyper}}{\usepackage{footnote}}
\makesavenoteenv{longtable}
\setlength{\emergencystretch}{3em} % prevent overfull lines
\providecommand{\tightlist}{%
  \setlength{\itemsep}{0pt}\setlength{\parskip}{0pt}}
\usepackage{bookmark}
\IfFileExists{xurl.sty}{\usepackage{xurl}}{} % add URL line breaks if available
\urlstyle{same}
% fallback for those not using the hyperref driver hyperxmp:
\makeatletter
\@ifundefined{xmpquote}{\newcommand{\xmpquote}[1]{#1}}{}
\makeatother
\hypersetup{
  colorlinks=true,
  linkcolor={Maroon},
  filecolor={Maroon},
  citecolor={Blue},
  urlcolor={Blue},
  pdfcreator={LaTeX via pandoc}}

\author{}
\date{}

\usepackage{pdflscape,fvextra}
\setlength{\tabcolsep}{3pt}
\begin{document}

\section{B2 in the published RECeSS TRANSCRIPT
runner}\label{b2-in-the-published-recess-transcript-runner}

\subsection{Exact comparison target}\label{exact-comparison-target}

This extension pins
\href{https://github.com/RECeSS-EU-Project/benchmark-code}{benchmark-code}
at commit \texttt{a7f11077271ced\allowbreak{}f3a98a82e3dc74\allowbreak{}b6fc0e93986e} and the
authors\textquotesingle{}
\href{https://github.com/RECeSS-EU-Project/benchmark-results}{benchmark-results}
at commit \texttt{cf5d9fdcb1ccd1\allowbreak{}676c7a0e7a39e7\allowbreak{}84d79557fae0}. The
comparison is limited to TRANSCRIPT v2.0.0 and its 11 published
reference models. B2 is registered in the upstream pipeline by
\texttt{benchmarks/\allowbreak{}recess\_\allowbreak{}adapter/\allowbreak{}official\_\allowbreak{}b2.\allowbreak{}patch}. The patch also
replaces Unix-only directory and cleanup commands with Python calls,
converts the generated NumPy seed to a Python integer, and restores an
alias required by the pinned \texttt{cute-ranking} package on NumPy 2.
These changes do not alter split, training, selection, or metric
functions. The run retained its small per-seed intermediate CSVs in the
ignored artifact directory; the versioned aggregate and seed CSVs are
the comparison inputs. The B2 adapter caches feature-only computations
across folds; its scores are tested for exact equality with the frozen
B2 implementation, including after validation labels are changed.

The upstream runner generates the same 100 seeds for each model with
\texttt{np.\allowbreak{}random.\allowbreak{}seed(1234)} followed by
\texttt{np.\allowbreak{}random.\allowbreak{}choice(range(int(1e8)),\ size=100)}. Each seed makes a
20\% outer split with \texttt{random\_\allowbreak{}simple\_\allowbreak{}split} or
\texttt{weakly\_\allowbreak{}correlated\_\allowbreak{}split}. \texttt{cv\_\allowbreak{}training} trains five
inner-fold models using \texttt{cv\_\allowbreak{}type="random"} and chooses the one
with the highest inner row-wise AUC. For B2, \texttt{neighbors=10} and
\texttt{rrf\_\allowbreak{}k=60} are fixed before these runs. The published
\texttt{K=5} operation is fold-based model selection; it does not search
a hyperparameter grid. Our separate three-fold B2 grid study remains a
distinct experiment and must not be mixed into this table. All 22
published TRANSCRIPT parameter JSON files (11 models × two split types)
set \texttt{params} to \texttt{null}, confirming that the reference
result files do not represent a five-fold hyperparameter grid search
either.

In \texttt{stanscofi} 2.0.1, \texttt{weakly\_\allowbreak{}correlated\_\allowbreak{}split} seeds
random-number generators but its clustering and fold assignment contain
no subsequent random draw. Consequently its 100 seeds repeat the same
outer holdout; only the inner five-fold split can change with the seed.
This is part of the published protocol, and the weak-split spread must
not be interpreted as uncertainty over 100 distinct drug-group holdouts.

The primary result is the upstream \texttt{rowwise\_\allowbreak{}metrics.\allowbreak{}calc\_\allowbreak{}auc}
output labeled \texttt{Lin\textquotesingle{}s\ AUC} in the CSV and
\texttt{NS\ AUC} in the authors\textquotesingle{} analysis. The same
runner also writes global AUC, global NDCG, row-wise AUC/NDCG and other
metrics. The summary script requires exact agreement of all 100 seed
positions for B2 and every published model before producing any paired
difference. It records source-file SHA-256 values and finite-value
coverage for each metric.

\subsection{Result}\label{result}

B2 achieves NS-AUC \texttt{0.\allowbreak{}5222\ ±\ 0.\allowbreak{}0354} on random simple (rank
9/12) and \texttt{0.\allowbreak{}5019\ ±\ 0.\allowbreak{}0036} on weakly correlated (rank 8/12).
The leading published references are BNNR (\texttt{0.\allowbreak{}7331}) and MBiRW
(\texttt{0.\allowbreak{}7384}), respectively. B2 is below both on all 100 paired
seeds. The \href{benchmark_results.md}{full table} and
\href{../benchmark/results/recess_official_b2_vs_11.json}{machine-readable
comparisons} include all 11 references, paired differences, and source
hashes. B2\textquotesingle s aggregate result, seed, and parameter CSVs
are versioned in \texttt{benchmark/\allowbreak{}results/\allowbreak{}recess\_\allowbreak{}official\_\allowbreak{}b2/\allowbreak{}}.

\subsection{Reproduce}\label{reproduce}

From the repository root in PowerShell:

\begin{Shaded}
\begin{Highlighting}[]
\VariableTok{$upstream} \OperatorTok{=} \FunctionTok{Join{-}Path} \VariableTok{$env}\OperatorTok{:}\VariableTok{TEMP} \StringTok{\textquotesingle{}recess{-}benchmark{-}code{-}pinned\textquotesingle{}}
\VariableTok{$published} \OperatorTok{=} \FunctionTok{Join{-}Path} \VariableTok{$env}\OperatorTok{:}\VariableTok{TEMP} \StringTok{\textquotesingle{}recess{-}benchmark{-}results{-}pinned\textquotesingle{}}
\NormalTok{git clone https}\OperatorTok{://}\NormalTok{github}\OperatorTok{.}\FunctionTok{com}\OperatorTok{/}\NormalTok{RECeSS{-}EU{-}Project}\OperatorTok{/}\NormalTok{benchmark{-}code}\OperatorTok{.}\FunctionTok{git} \VariableTok{$upstream}
\NormalTok{git }\OperatorTok{{-}}\NormalTok{C }\VariableTok{$upstream}\NormalTok{ checkout a7f11077271cedf3a98a82e3dc74b6fc0e93986e}
\VariableTok{$patch} \OperatorTok{=} \OperatorTok{(}\FunctionTok{Resolve{-}Path}\NormalTok{ benchmarks}\OperatorTok{/}\NormalTok{recess\_adapter}\OperatorTok{/}\NormalTok{official\_b2}\OperatorTok{.}\FunctionTok{patch}\OperatorTok{).}\FunctionTok{Path}
\NormalTok{git }\OperatorTok{{-}}\NormalTok{C }\VariableTok{$upstream}\NormalTok{ apply }\OperatorTok{{-}{-}}\NormalTok{recount }\VariableTok{$patch}
\NormalTok{git clone }\OperatorTok{{-}{-}}\NormalTok{filter}\OperatorTok{=}\NormalTok{blob}\OperatorTok{:}\NormalTok{none }\OperatorTok{{-}{-}}\NormalTok{sparse https}\OperatorTok{://}\NormalTok{github}\OperatorTok{.}\FunctionTok{com}\OperatorTok{/}\NormalTok{RECeSS{-}EU{-}Project}\OperatorTok{/}\NormalTok{benchmark{-}results}\OperatorTok{.}\FunctionTok{git} \VariableTok{$published}
\NormalTok{git }\OperatorTok{{-}}\NormalTok{C }\VariableTok{$published}\NormalTok{ checkout cf5d9fdcb1ccd1676c7a0e7a39e784d79557fae0}
\NormalTok{git }\OperatorTok{{-}}\NormalTok{C }\VariableTok{$published}\NormalTok{ sparse{-}checkout }\FunctionTok{set}\NormalTok{ results\_TRANSCRIPT results\_TRANSCRIPT\_weakly\_correlated}
\NormalTok{python }\OperatorTok{{-}}\NormalTok{m scripts}\OperatorTok{.}\FunctionTok{run\_recess\_official\_b2} \OperatorTok{{-}{-}}\NormalTok{upstream }\VariableTok{$upstream} \OperatorTok{{-}{-}}\NormalTok{n }\DecValTok{100} \OperatorTok{{-}{-}}\NormalTok{k }\DecValTok{5}
\DataTypeTok{python} \OperatorTok{{-}}\NormalTok{m scripts}\OperatorTok{.}\FunctionTok{compare\_recess\_official\_b2} \OperatorTok{{-}{-}}\NormalTok{published }\VariableTok{$published}
\CommentTok{\# To recompute only the comparison from the committed B2 CSVs:}
\NormalTok{python }\OperatorTok{{-}}\NormalTok{m scripts}\OperatorTok{.}\FunctionTok{compare\_recess\_official\_b2} \OperatorTok{{-}{-}}\NormalTok{published }\VariableTok{$published} \OperatorTok{{-}{-}}\NormalTok{ours benchmark}\OperatorTok{/}\NormalTok{results}\OperatorTok{/}\NormalTok{recess\_official\_b2}
\end{Highlighting}
\end{Shaded}

Install the environment dependencies from
\texttt{requirements-benchmark.\allowbreak{}lock} and \texttt{statsmodels} before
running. The runner stages three checksum-verified input CSV files from
\texttt{data/\allowbreak{}raw/\allowbreak{}TRANSCRIPT\_\allowbreak{}dataset\_\allowbreak{}v2.\allowbreak{}0.\allowbreak{}0} using hard links when
possible. Its default outputs are under ignored
\texttt{artifacts/\allowbreak{}recess\_\allowbreak{}official\_\allowbreak{}b2}; the comparison summary is
\texttt{benchmark/\allowbreak{}results/\allowbreak{}recess\_\allowbreak{}official\_\allowbreak{}b2\_\allowbreak{}vs\_\allowbreak{}11.\allowbreak{}json}. For a
smoke test, use \texttt{-\/\allowbreak{}-n\ 1\ -\/\allowbreak{}-splitting\ random\_\allowbreak{}simple} and a
separate \texttt{-\/\allowbreak{}-output} directory. The upstream runner reuses a
finished result CSV, so do not reuse a smoke-test directory for the
100-run comparison.

\subsection{Interpretation boundary}\label{interpretation-boundary}

Matching seeds, folds, source data version, and metric code makes this a
protocol-matched comparison against the 11 author-published result
files. The 11 models are not rerun in the local environment. Their
original raw dataset bytes and software builds are not independently
checksum verified by the published CSVs, so any environment or
source-data discrepancy remains a reproducibility limitation. The
primary metric is a ranking measure for known and unknown associations;
it is not a clinical benefit estimate.

As a cross-check of the local patched runner and TRANSCRIPT input, an
independent \texttt{LogisticMF} run on the first official random seed
(\texttt{17263407}) exactly reproduced the authors\textquotesingle{}
published values for NS-AUC (\texttt{0.\allowbreak{}6704512218407673}), global AUC
(\texttt{0.\allowbreak{}950127861336678}), and global NDCG
(\texttt{0.\allowbreak{}5501272684100554}). This checks one reference run, not all 11
× 100 runs.

\subsection{B3: row-oriented fusion
(2026-09-28)}\label{b3-row-oriented-fusion-2026-09-28}

The official NS-AUC ("Lin\textquotesingle s AUC",
\texttt{rowwise\_\allowbreak{}metrics.\allowbreak{}calc\_\allowbreak{}auc(.\allowbreak{}.\allowbreak{}.\allowbreak{},\ transpose=False)}) ranks
diseases \textbf{within each drug row} and counts tied scores as losses
(strict \texttt{\textgreater{}}). B2 ranked candidates within disease
columns, and its popularity component is constant within a drug row. B3
(\texttt{benchmarks/\allowbreak{}recess\_\allowbreak{}adapter/\allowbreak{}official\_\allowbreak{}b3.\allowbreak{}py}) rank-normalises
six training-free components within drug rows and averages them: disease
popularity, drug- and disease-expression kNN, label co-occurrence
drug/disease CF, and reversal. It adds a deterministic
disease-popularity tie-breaker.

The configuration was frozen in
\texttt{configs/\allowbreak{}b3\_\allowbreak{}row\_\allowbreak{}fusion\_\allowbreak{}v1.\allowbreak{}json} at commit \texttt{92943bb},
before any official scoring. It was selected on eight random-split
development seeds disjoint from the 100 official seeds.
\texttt{weakly\_\allowbreak{}correlated\_\allowbreak{}split} has no seed-dependent draw, so its
development run necessarily used the official outer holdout; this is
disclosed in the config and was not used for selection.

{\def\LTcaptype{none} % do not increment counter
\begin{longtable}[]{@{}>{\raggedright\arraybackslash}p{0.1800\textwidth}>{\raggedright\arraybackslash}p{0.1800\textwidth}>{\raggedright\arraybackslash}p{0.1800\textwidth}>{\raggedright\arraybackslash}p{0.1800\textwidth}>{\raggedright\arraybackslash}p{0.1800\textwidth}@{}}
\toprule\noalign{}
Split & B3 NS-AUC (mean ± SD, 100 seeds) & Rank / 13 & Best published &
B3 paired wins vs best \\
\midrule\noalign{}
\endhead
\bottomrule\noalign{}
\endlastfoot
Random simple & 0.7234 ± 0.0339 & 2 & BNNR 0.7331 & 35/100 \\
Weakly correlated & 0.6919 ± 0.0088 & 3 & MBiRW 0.7384 & 0/100 \\
\end{longtable}
}

B3 does not surpass the best published model on either split. The gain
from B2 (0.5222 → 0.7234) comes from matching the
metric\textquotesingle s orientation and from known-association
structure, not from signature reversal. Machine-readable comparison:
\texttt{benchmark/\allowbreak{}results/\allowbreak{}recess\_\allowbreak{}official\_\allowbreak{}b3\_\allowbreak{}vs\_\allowbreak{}11.\allowbreak{}json}; aggregate
CSVs: \texttt{benchmark/\allowbreak{}results/\allowbreak{}recess\_\allowbreak{}official\_\allowbreak{}b3/\allowbreak{}}. The post-hoc
variant B3noREV (reversal removed) completed afterwards: 0.7124 random
simple (rank 2/13) and 0.7050 weakly correlated (rank 2/13, behind
MBiRW). Because REV removal was chosen after the weak-split development
score had been seen, B3noREV is reported only as a secondary, post-hoc
result; B3 remains the primary pre-registered model. Ranks here are
within \{11 published models, B2, the model\}
(\texttt{rank\_\allowbreak{}in\_\allowbreak{}published\_\allowbreak{}field}).

\subsection{B4: round 2, B3 plus a NumPy BNNR port
(2026-09-29)}\label{b4-round-2-b3-plus-a-numpy-bnnr-port-2026-09-29}

BNNR\textquotesingle s published implementation calls Octave.
\texttt{benchmarks/\allowbreak{}recess\_\allowbreak{}adapter/\allowbreak{}bnnr\_\allowbreak{}numpy.\allowbreak{}py} ports \texttt{BNNR.\allowbreak{}m}
and \texttt{svt.\allowbreak{}m} line by line. On official random-split seeds 1--3 it
reproduces the published per-seed NS-AUC with a maximum absolute
difference of 0.0
(\texttt{benchmark/\allowbreak{}results/\allowbreak{}bnnr\_\allowbreak{}numpy\_\allowbreak{}fidelity.\allowbreak{}json}). Those seeds
were used only to check fidelity.

B4 averages within-drug-row ranks of B3 and BNNR at weights 1:2. It was
chosen among five pre-listed ensembles on 10 random-split development
seeds disjoint from the official 100, and frozen at commit
\texttt{e5ff7e7} (\texttt{configs/\allowbreak{}b4\_\allowbreak{}ensemble\_\allowbreak{}v1.\allowbreak{}json}) before any
official B4 score existed. The official B3 results were already public
at that point, so B4 is a second, separately labelled round.

\textbf{Split Random simple · 0.7453 ± 0.0333}

\begin{itemize}
\tightlist
\item
  \textbf{Rank / 13：} 1
\item
  \textbf{Best published：} BNNR 0.7331
\item
  \textbf{B4 paired wins vs best：} 87/100
\item
  \textbf{Mean paired difference：} +0.0122
\end{itemize}

\textbf{Split Weakly correlated · 0.6585 ± 0.0187}

\begin{itemize}
\tightlist
\item
  \textbf{Rank / 13：} 3
\item
  \textbf{Best published：} MBiRW 0.7384
\item
  \textbf{B4 paired wins vs best：} 0/100
\item
  \textbf{Mean paired difference：} -0.0799
\end{itemize}

B4 ranks first on random simple but not on weakly correlated, where
MBiRW (0.7384) leads. It is therefore not state of the art across both
protocols. The weakly correlated protocol repeats one outer holdout
across all seeds, so its paired values are not independent replications.
B4 still trains nothing new: BNNR is a published matrix-completion
method run through the official protocol, and the ensemble weights were
fixed on development seeds. Comparison:
\texttt{benchmark/\allowbreak{}results/\allowbreak{}recess\_\allowbreak{}official\_\allowbreak{}b4\_\allowbreak{}vs\_\allowbreak{}11.\allowbreak{}json}. Correction
(2026-10-01): seed test sets overlap, and under the Nadeau--Bengio
corrected test the B4 − BNNR difference (+0.0122) has p = 0.066, 95\% CI
{[}-0.0008, +0.0251{]}. The random-split lead is therefore not
significant, and BNNR remains stronger on global NDCG and HR@10
(\href{benchmark_stats_v1.md}{benchmark\_stats\_v1}).

\end{document}
